\begin{figure*}
\centering
\includegraphics[width=\textwidth]{figs/appendix-trace-mlp-detail-crop.pdf}
\caption{\small Detail view of causal traces, breaking out a representative set of individual cases from the 1000 factual statements that are averaged in \reffig{trace-mlp-disabled}.  Shows the causal trace at a specific subject token, with and without MLP disabled, as described in \refsec{causal-tracing}.  In every case, the token tested is highlighted in a red box.  In (a,b,c,d,e) cases are shown that fit the typical pattern: Restoring individual hidden states at a range of layers has a strong decisive average causal effect at the last token of the subject.  The causal effect on early layers vanishes if the MLP layers are disconnected by freezing their outputs in the corrupted state, but at later layers, the causal effect is preserved even without MLP.   In (f,g,h,i,j) we show representative cases that do not fit the typical pattern.  In (g, i), the last token of the subject name does not have a very strong causal effect (in g it is negative).  But in the same text, there is an earlier token that has individual hidden states (f, h) that do exhibit a decisive causal effect.  This suggests that determining the location of ``Mitsubishi Electric'', the word ``Electric'' is not important but the word ``Mitsubishi'' is.  Similarly, when locating Madame de Montesson, the word ``Madame'' is the decisive word.  (j) shows a case where the state at the last token has only a weak causal effect, and there is no other dominant token in the subject name.
}
\lblfig{apd-trace-mlp-detail}
\end{figure*}